% ---- pg-345 （书页 333） ----

\noindent 3.\quad Atkinson, A. C. \textit{Plots, Transformations and Regression: An Introduction to Graphical Methods of Diagnostic Regression Analysis}. Oxford, England: Oxford University Press, 1986.

\noindent 4.\quad Bartlett, M. S. “A Note on Multiplying Factors for Various Chi-Squared Approximations.” \textit{Journal of the Royal Statistical Society (B)}, \textbf{16} (1954), 296–298.

\noindent 5.\quad Belsley, D. A., E. Kuh, and R. E. Welsh. \textit{Regression Diagnostics: Identifying Influential Data and Sources of Collinearity} (Paperback). New York: Wiley-Interscience, 2004.

\noindent 6.\quad Bowerman, B. L., and R. T. O'Connell. \textit{Linear Statistical Models: An Applied Approach} (2nd ed.). Belmont, CA: Thompson Brooks/Cole, 2000.

\noindent 7.\quad Box, G. E. P. “A General Distribution Theory for a Class of Likelihood Criteria.” \textit{Biometrika}, \textbf{36} (1949), 317–346.

\noindent 8.\quad Box, G. E. P., G. M. Jenkins, and G. C. Reinsel. \textit{Time Series Analysis: Forecasting and Control} (3rd ed.). Englewood Cliffs, NJ: Prentice Hall, 1994.

\noindent 9.\quad Chatterjee, S., A. S. Hadi, and B. Price. \textit{Regression Analysis by Example} (4th ed.). New York: Wiley-Interscience, 2006.

\noindent 10.\quad Cook, R. D., and S. Weisberg. \textit{Applied Regression Including Computing and Graphics}. New York: John Wiley, 1999.

\noindent 11.\quad Cook, R. D., and S. Weisberg. \textit{Residuals and Influence in Regression}. London: Chapman and Hall, 1982.

\noindent 12.\quad Daniel, C. and F. S. Wood. \textit{Fitting Equations to Data} (2nd ed.) (paperback). New York: Wiley-Interscience, 1999.

\noindent 13.\quad Draper, N. R., and H. Smith. \textit{Applied Regression Analysis} (3rd ed.). New York: John Wiley, 1998.

\noindent 14.\quad Durbin, J., and G. S. Watson. “Testing for Serial Correlation in Least Squares Regression, II.” \textit{Biometrika}, \textbf{38} (1951), 159–178.

\noindent 15.\quad Galton, F. “Regression Toward Mediocrity in Heredity Stature.” \textit{Journal of the Anthropological Institute}, \textbf{15} (1885), 246–263.

\noindent 16.\quad Goldberger, A. S. \textit{Econometric Theory}. New York: John Wiley, 1964.

\noindent 17.\quad Heck, D. L. “Charts of Some Upper Percentage Points of the Distribution of the Largest Characteristic Root.” \textit{Annals of Mathematical Statistics}, \textbf{31} (1960), 625–642.

\noindent 18.\quad Khattree, R. and D. N. Naik. \textit{Applied Multivariate Statistics with SAS\textsuperscript{\textregistered} Software} (2nd ed.) Cary, NC: SAS Institute Inc., 1999.

\noindent 19.\quad Lee, J. “Relationships Between Properties of Pulp-Fibre and Paper.” Unpublished doctoral thesis, University of Toronto, Faculty of Forestry, 1992.

\noindent 20.\quad Neter, J., W. Wasserman, M. Kutner, and C. Nachtsheim. \textit{Applied Linear Regression Models} (3rd ed.). Chicago: Richard D. Irwin, 1996.

\noindent 21.\quad Pillai, K. C. S. “Upper Percentage Points of the Largest Root of a Matrix in Multivariate Analysis.” \textit{Biometrika}, \textbf{54} (1967), 189–193.

\noindent 22.\quad Rao, C. R. \textit{Linear Statistical Inference and Its Applications} (2nd ed.) (paperback). New York: Wiley-Interscience, 2002.

\noindent 23.\quad Seber, G. A. F. \textit{Linear Regression Analysis}. New York: John Wiley, 1977.

\noindent 24.\quad Rudorfer, M. V. “Cardiovascular Changes and Plasma Drug Levels after Amitriptyline Overdose.” \textit{Journal of Toxicology-Clinical Toxicology}, \textbf{19} (1982), 67–71.
